\documentclass[10pt,a4paper]{article}
\usepackage[T1]{fontenc}
\usepackage[utf8]{inputenc}
\usepackage[english]{babel}
\usepackage{lmodern}
\usepackage[margin=22mm]{geometry}
\usepackage{booktabs,tabularx,hyperref}
\setlength{\parindent}{0pt}
\setlength{\parskip}{5pt}
\title{File Streaming and Unicode Tests\\\large Raspberry Pi 3 -- Team C}
\author{Guillem Moreno}
\date{October 2026}
\begin{document}
\maketitle
\section{Scope}
This directory provides equivalent Python 3.10+ and C++17 implementations for
reading a file as bytes, supplying a transport stream, reconstructing the original
file, identifying Unicode encoding candidates and producing a separate UTF-8 text
copy. Both endpoints may run Raspberry Pi OS on a Pi 3.

The transmitted stream is exactly the source file, including any BOM, NUL bytes
and original line endings. It contains no new header, filename or encoding tag.
The caller supplies the original byte count and CRC-32 separately. These values
are application metadata for this test harness, not a proposed radio frame format.
CRC-32 detects accidental corruption; it is not authentication.

The implementation uses bounded buffers: 4,096 bytes by default, configurable from
1 byte to 1 MiB. Inspection makes several streaming validation passes rather than
loading the whole file. Source files must remain unchanged during inspection and
transfer. A size/CRC check detects changes between preparation and transmission.

\section{Build and run on Raspberry Pi OS}
Run from the repository root:
\begin{verbatim}
sudo apt-get install build-essential python3
make -C proto/test_encoding_text test
cd proto/test_encoding_text
python3 generate_samples.py build/samples
python3 encoding_stream.py inspect build/samples/utf-8_bom.txt
./build/encoding_stream inspect build/samples/utf-8_bom.txt
\end{verbatim}
The Python implementation uses only the standard library. The C++ implementation
has no external library dependencies. Build it with a native Pi compiler;
the supplied CI also builds ARM32 hard-float and AArch64 executables. Cross-builds
are compile checks, not evidence of execution on a Pi or a physical USB stick.

Both command-line programs expose the same commands and emit JSON results.
Use \texttt{--help} for arguments. The examples below use the Python script.
The C++ executable is \texttt{build/encoding\_stream}.

\begin{verbatim}
python3 encoding_stream.py discover /media/pi/INPUT
python3 encoding_stream.py copy input.txt recovered.txt
python3 encoding_stream.py send input.txt stream.bin
\end{verbatim}
Discovery accepts exactly one regular \texttt{.txt} file in the specified directory
(case-insensitive extension). It ignores directories and symlinks, and rejects
zero or multiple candidates. An explicit source path bypasses discovery.

\texttt{send} prints the byte count and CRC-32. Pass those exact values to
\texttt{receive}; for example, the nine bytes \texttt{123456789} have length 9 and
CRC-32 \texttt{cbf43926}:
\begin{verbatim}
python3 encoding_stream.py receive stream.bin recovered.txt \
    --size 9 --crc32 cbf43926
\end{verbatim}
Use \texttt{-} as the stream path to read standard input or write standard output.
When sending to stdout, metadata is printed to stderr so it cannot contaminate
the byte stream. The stream may be moved between machines by the existing
transport; a pipe or stream file is sufficient for local testing.

\section{Encoding selection}
The supported encodings are UTF-8, UTF-16 LE/BE and UTF-32 LE/BE. The operating
system reads bytes; the application chooses how those bytes are interpreted.
The inspection result includes all matching BOM signatures, all syntactically
valid encodings, the automatic selection if available, and a detection status.

\begin{tabularx}{\linewidth}{@{}lX@{}}
\toprule
Status & Meaning\\
\midrule
\texttt{bom} & One valid BOM interpretation.\\
\texttt{ambiguous\_bom} & Multiple valid interpretations of an initial signature.\\
\texttt{invalid\_bom} & Recognised signature, but no valid matching text.\\
\texttt{unique\_valid} & Exactly one valid candidate within the five supported encodings.\\
\texttt{ambiguous} & No BOM and several valid candidates.\\
\texttt{invalid} & No supported encoding validates.\\
\texttt{empty} & No bytes from which to infer an encoding.\\
\bottomrule
\end{tabularx}

Automatic selection is conservative. Missing BOMs do not imply UTF-8.
For example, \texttt{41 00} can be UTF-8 A followed by NUL or UTF-16 LE A.
Similarly, \texttt{FF FE 00 00 41 00 00 00} permits a UTF-32 LE BOM plus A,
or a UTF-16 LE BOM plus NUL, A, NUL. Strict automatic decoding rejects this
overlap. Use an explicitly declared encoding to resolve it:
\begin{verbatim}
python3 encoding_stream.py decode recovered.txt display.txt \
    --encoding utf-16-le --chunk-size 3
\end{verbatim}
The optional \texttt{--bom-policy longest} applies a conventional longest-signature
rule. It is suitable only when that convention is agreed for the input files;
it does not prove the source encoding. A unique valid candidate without a BOM is
also a conclusion limited to the supported set, not a universal detector.

Decoding validates strictly: malformed sequences, truncated code units, isolated
surrogates, overlong UTF-8 and values above U+10FFFF cause an error. NUL and Unicode
noncharacters remain valid. The decoder maintains state across chunks, including
split UTF-8 sequences and UTF-16 surrogate pairs.

The decoded output is a separate UTF-8 file. One initial BOM is consumed by
default; \texttt{--keep-bom} retains the initial U+FEFF as text. Interior U+FEFF,
combining sequences and line endings are preserved. The original file is never
normalised, re-encoded or overwritten.

\section{Receiver and filesystem behaviour}
\texttt{FileReceiver} accepts ordered data with byte offsets. Gaps, duplicate
offsets, out-of-order chunks and excess data are rejected. The transport must
perform ordering, deduplication and retries before delivering bytes to this layer.
Finalisation checks the expected length and CRC-32, flushes and syncs the file,
then reads it back for verification.

Data is first written to \texttt{destination.part}, created exclusively.
An existing final file or partial file is never overwritten. On success, Linux
\texttt{renameat2(RENAME\_NOREPLACE)} publishes the result, and the parent directory
is synced. Raspberry Pi OS must provide this operation on the destination
filesystem. Unsupported publication or sync operations are reported as errors.

A normal exception removes the partial file owned by that operation. A killed
process or power loss can leave a partial file; inspect it and remove it explicitly
before retrying. This test layer does not implement resume. If directory sync fails
after publication, the complete final file can exist although the operation
reports failure. It must be inspected before a retry. Safely unmount USB media
after files are closed.

Incomplete transfers are not published as complete files. The existing competition
node can save a received prefix at STOP; that policy is separate from this
complete-file receiver and must not be silently replaced by it.

\section{API integration}
Python exposes \texttt{file\_info}, \texttt{file\_chunks}, \texttt{send\_file},
\texttt{FileReceiver}, \texttt{inspect\_file} and \texttt{decode\_file}.
C++ exposes equivalent functions in namespace \texttt{mtp\_text}.
\begin{verbatim}
# Python: ordered payload delivery by the existing transport.
info = file_info(source_path)
with FileReceiver(destination_path, info) as receiver:
    for offset, payload in transport_payloads:
        receiver.accept(offset, payload)
    receiver.finish()
\end{verbatim}
Here \texttt{transport\_payloads} represents the transport adapter, not a function
implemented by this test package. The C++ sender takes a byte-sink callback; the
receiver takes \texttt{accept(offset, data, length)}. Both expose a finalisation
call that the adapter invokes only when transfer completion has been established.

The repository already supplies \texttt{proto/src/codec.c} and
\texttt{proto/port/linux/store.c}. When using its compression path, compression
operates on the original byte stream and decompression precedes this receiver's
file checks. Existing radio frames, retransmission state and compressed-stream
metadata are unchanged. This package is a test harness and reusable API; it is
not wired into \texttt{teamc-node} automatically.

This harness uses files on a mounted filesystem. The existing node accesses its
USB stick through FatFs as a raw block device. Do not use both paths concurrently
on the same mounted device. Integration must use one ownership model consistently.

\section{Validation}
\texttt{make test} builds the C++ program and API checks, then runs Python unit
tests and cross-language tests. Cases include all five encodings with and without
BOMs; empty and BOM-only files; NUL, combining marks, emoji and mixed line endings;
overlapping and incomplete signatures; malformed sequences; short writes;
truncation, corruption and excess data; existing destinations and publication
races; offsets and file discovery. Differential tests compare both inspectors on
fixed and deterministic random inputs. Streaming tests cover multi-megabyte files.

Run the tests on both Pis and on the actual destination filesystem before hardware
integration. Host tests and ARM compilation do not verify radio timing, physical
USB behaviour or resilience to a real power interruption.
\end{document}
